\subsection{Vertex cover from a maximal matching}
\label{gre:vc}
View each edge as an element and each vertex as the set of its incident edges. Every
element then belongs to exactly two sets. For the \emph{unweighted} problem, greedily
construct any maximal matching $M$ and return the set $C$ of \textbf{all endpoints of $M$}.
\begin{algorithm}[H]
\caption{Maximal-matching vertex cover}
\KwIn{Undirected graph $G=(V,E)$}
\KwOut{A vertex cover $C\subseteq V$}
$M\gets\varnothing$; mark all vertices unmatched\;
\ForEach{$\{u,v\}\in E$}{
  \If{both $u$ and $v$ are unmatched}{
    $M\gets M\cup\{\{u,v\}\}$; mark $u,v$ matched\;
  }
}
\Return $C=\bigcup_{e\in M}e$\;
\end{algorithm}
\begin{proof}
An uncovered edge would have two unmatched endpoints and could be added to $M$, contrary
to maximality. Any vertex cover must take a distinct endpoint from each pairwise disjoint
matching edge, so $|M|\le\OPT$. Hence $|C|=2|M|\le2\OPT$.
The scan takes $O(|V|+|E|)$ time. Maximum matching is unnecessary; maximality proves
coverage, while disjointness proves the lower bound.
\end{proof}
This unweighted cardinality proof does not give a weighted-vertex-cover guarantee by
simply summing endpoint weights; use LP rounding for that version.

\paragraph{Why complements do not preserve approximation.}
The complement of any vertex cover is an independent set, and $\OPT_{IS}=|V|-\OPT_{VC}$.
But $|C|\le2\OPT_{VC}$ only implies
$|V\setminus C|\ge2\OPT_{IS}-|V|$, which can be useless. On $q$ disjoint edges plus one
isolated vertex, the matching algorithm selects all $2q$ nonisolated vertices. Its complement
has size 1, while $\OPT_{IS}=q+1$. Thus exact optimization equivalence does not imply
approximation preservation.

\subsection{Triangle deletion and disjoint-obstruction packing}
\label{gre:triangles}
To delete as few edges as possible so that an undirected graph becomes triangle-free,
repeatedly find a remaining triangle and delete \emph{all three} of its edges. Return the
union $F$ of deleted edges.
\begin{proof}
The process stops with no triangle. The selected triangles are pairwise edge-disjoint,
because earlier selected edges were deleted. If $k$ triangles were selected, every feasible
deletion set must remove at least one distinct edge from each of these triangles, so
$\OPT\ge k$. Our output has $|F|=3k\le3\OPT$.
One implementation scans all vertex triples once with an adjacency matrix: select a
triangle only if none of its edges was selected earlier. At the end, any surviving triangle
would have been eligible when its triple was scanned, a contradiction. Time is
$O(|V|^3+|E|)$ with constant-time edge tests.
\end{proof}
\textbf{Counterexample to deleting an arbitrary single edge per triangle:}
take a shared edge $ab$ and vertices $v_1,\ldots,v_q$, each joined to both $a$ and $b$.
Deleting $ab$ solves the instance, while repeatedly deleting a spoke can require $q$ edges.

\paragraph{Reusable generalization.}
For fixed forbidden subgraphs whose occurrences cannot be created by deleting edges,
suppose every forbidden configuration has $r$ edges for a fixed constant $r$. Select a maximal
collection of edge-disjoint configurations and delete all their edges. Maximality gives
feasibility, disjointness gives $\OPT\ge k$, and $\ALG=rk\le r\OPT$. Detection must be
polynomial; the fact that $r$ is fixed is relevant. Induced patterns need separate care:
edge deletion can create a new induced pattern. Weighted deletions need a different cost argument.

\paragraph{Greedy by number of triangles removed.}
This is weighted set cover on the universe of original triangles, with each graph edge
covering the triangles that contain it. Recompute marginal coverage after each deletion.
The generic proof gives an $H_N$ bound, where $N$ is the number of original triangles;
it does not by itself establish the factor 3 above.
